\subsection{N-010: an explicit split form and GT path basis over $\mathbb Q(D)$}
\label{sec:nonstandard-rational-gt}

\paragraph{Statement and scope.}
Let $F=\mathbb Q(D)$, $D=(v+v^{-1})^2$, and let
$\mathcal B_{r,F}$ be the $F$-span of the actual $P$-words in the
generic nonstandard algebra.  The fact that this is an $F$-form is the
supplied counting consult, Theorem 1.  The full algebra at rank four is
split over $F$:
\[
 \mathcal B_{4,F}\cong
 F\times M_2(F)\times M_3(F)\times M_6(F)
 \times M_5(F)\times M_5(F)\times M_2(F)\times M_3(F)\times F.
 \tag{N-25}
\]
The tower through rank four has multiplicity-free branching, and every
GT path line, together with a concrete choice of vector on it, can be
constructed over $F$.  The accompanying certificate gives all nine
triples of generator matrices, all GT transition matrices, all path
labels, and an exact rank certificate for (N-25).

This does not identify these choices with VIII's $E(4)$.  It proves that
neither splitting the algebra nor choosing its GT path vectors forces
an algebraic extension of $\mathbb Q(D)$ at rank four.  The separate
minimum-degree globalization and its prescribed normalization remain
uncomputed.

\paragraph{Three small Hecke models.}
Put $\delta=v+v^{-1}$.  The following matrices represent
$A_i=\delta C'_{s_i}$, with column-vector convention:
\[
 A_i^{(4)}=(D)\qquad(i=1,2,3),
\]
\[
 A_1^{(31)}=\begin{pmatrix}0&0&0\\D&D&0\\0&0&D\end{pmatrix},\quad
 A_2^{(31)}=\begin{pmatrix}D&1&0\\0&0&0\\0&1&D\end{pmatrix},\quad
 A_3^{(31)}=\begin{pmatrix}D&0&0\\0&D&D\\0&0&0\end{pmatrix},
 \tag{N-26}
\]
\[
 A_1^{(22)}=A_3^{(22)}=
 \begin{pmatrix}D&1\\0&0\end{pmatrix},\qquad
 A_2^{(22)}=\begin{pmatrix}0&0\\D&D\end{pmatrix}.
 \tag{N-27}
\]
For each triple, direct multiplication verifies
\[
 A_i^2=DA_i,\qquad A_1A_3=A_3A_1,\qquad
 A_iA_{i+1}A_i-DA_i=A_{i+1}A_iA_{i+1}-DA_{i+1}.
 \tag{N-28}
\]
Consequently $T_i=\delta^{-1}A_i-v^{-1}I$ obey the ordinary
Hecke quadratic and braid relations, so these are actual Hecke
representations after extension to $\mathbb Q(v)$.  In the braid
calculation, putting $a=\delta v^{-1}$ gives
$a(a-D)=-D$, which proves the equivalence with the last relation
in (N-28).  The quadratic relation follows from $A_i^2=DA_i$.

For provenance, the script reconstructs (N-26)--(N-27) from the supplied
exact KL arithmetic.  On a $C'$-cell basis indexed by permutations $w$,
rescale by $\delta^{\ell(w)\bmod2}$.  A diagonal coefficient $\delta$
in $C'_i$ becomes $D$ in $A_i$; an off-diagonal integer coefficient
joining opposite parities becomes an integer or an integer times $D$.
The script checks these assertions coefficient by coefficient.  In this
$C'$ convention the RSK cell shape is the transpose of the displayed
Specht label: the identity cell is the sign representation.  More
fundamentally, the explicit checks (N-28) suffice for the construction,
independently of a general cell-module theorem.

\paragraph{Nine actual nonstandard modules over $F$.}
For two triples $A=(A_i)$ and $B=(B_i)$ above, the actual tensor
generator $P_i=C'_i\otimes C'_i+C_i\otimes C_i$ acts as
\[
 \mathsf P_i(A,B)=\frac{2}{D}A_i\otimes B_i
              -A_i\otimes I-I\otimes B_i+D I.
 \tag{N-29}
\]
Every entry is in $F$.  For $A=B$, (N-29) commutes with the tensor flip;
hence its symmetric and alternating subspaces are defined over $F$.
The symmetric square contains a one-dimensional simultaneous
$P_i=D$ eigenspace.  The script computes that line exactly over $F$
and takes the quotient, denoted $S'M_{\lambda}$; its dimensions are
two for $\lambda=(2,2)$ and five for $\lambda=(3,1)$.

Twisting by $\Theta(P_i)=D-P_i$ is also defined over $F$.
Indeed, the ordinary Hecke involution $T_i\mapsto-T_i^{-1}$ on one
tensor factor sends $C'_i\mapsto-C_i$ and $C_i\mapsto-C'_i$, so it
sends $P_i$ to $D-P_i$.  This proves validity of the twist in the
full ambient algebra, not merely in a quadratic presentation.

The nine constructed modules, in the order of (N-25), are
\[
 \epsilon_-,\ S'M_{22},\ \Lambda^2M_{31},\ M_{31}\otimes M_{22},\
 S'M_{31},\ (S'M_{31})^\Theta,\ M_4\otimes M_{22},\
 M_4\otimes M_{31},\ \epsilon_+.
 \tag{N-30}
\]
Here $\epsilon_-$ has $P_i=0$, and $\epsilon_+$ has $P_i=D$.
These names agree with the known list in GCT IV, subsection 11.8,
printed p.~67.  The construction itself provides a homomorphism
from $\mathcal B_{4,F}$ to the product of the nine endomorphism
algebras.  There is no circular assumption that these modules are
simple or pairwise distinct: that follows from the rank check next.

\paragraph{Exact splitting certificate.}
Flatten the nine matrix blocks of an element, in their displayed order,
to a vector with $1+4+9+36+25+25+4+9+1=114$ coordinates.  Specialize
$D=7$ and reduce modulo the prime $1000003$.  All denominators occurring
in the generator matrices are invertible there.  Exact Gaussian
elimination of actual word vectors gives rank $114$; the certificate
lists the 114 accepted words and every pivot column.  Its longest word
has length nine.  A determinant that is nonzero after a valid
specialization is nonzero in $F$, so the generic image has dimension
114.  The inherited generic dimension of the source is also 114.
The homomorphism is therefore an isomorphism, proving (N-25), absolute
simplicity of each factor module, and pairwise nonisomorphism.
This is a finite rank proof of the complete algebra decomposition,
not an inference that a short word search presents the integral order.

\paragraph{The exact branching projectors.}
Write $N_1,N_2$ for the two two-dimensional rank-three simples,
with $P_1P_2P_1=\lambda_iP_1$, where
$\lambda_1=D$ and $\lambda_2=(D-2)^2$.  In each of the nine models put
\[
 Z=(P_1-P_2)^2,\qquad z_0=0,\quad z_1=D(D-1),\quad z_2=4(D-1).
\]
The supplied exact arithmetic checks
$Z(Z-z_1I)(Z-z_2I)=0$ and that $Z$ commutes with $P_1,P_2$.
The distinct roots give the spectral idempotents
\[
 E_j=\prod_{k\ne j}\frac{Z-z_kI}{z_j-z_k}\quad(0\le j\le2).
 \tag{N-31}
\]
The exact checks also verify $E_0(P_1-P_2)=0$ and
$E_j(P_1P_2P_1-(D^2-z_j)P_1)=0$ for $j=1,2$.
On the zero eigenspace, $E_0P_1/D$ and $E_0(I-P_1/D)$ are the
$\epsilon_+$ and $\epsilon_-$ projectors.  The other two are the
$N_1$ and $N_2$ projectors.  Their exact ranks yield:
\[
\begin{array}{c|l}
\mathcal B_{4,F}\text{-module}&\text{restriction to }\mathcal B_{3,F}\\\hline
\epsilon_-&\epsilon_-\\
S'M_{22}&N_2\\
\Lambda^2M_{31}&\epsilon_-\oplus N_1\\
M_{31}\otimes M_{22}&\epsilon_+\oplus\epsilon_-\oplus N_1\oplus N_2\\
S'M_{31}&\epsilon_+\oplus N_1\oplus N_2\\
(S'M_{31})^\Theta&\epsilon_-\oplus N_1\oplus N_2\\
M_4\otimes M_{22}&N_1\\
M_4\otimes M_{31}&\epsilon_+\oplus N_1\\
\epsilon_+&\epsilon_+
\end{array}
 \tag{N-32}
\]
Each constituent occurs once.  For rank three to rank two, each $N_i$
has one $P_1=D$ and one $P_1=0$ line.  Therefore the nonzero projectors
$E_jP_1/D$ and $E_j(I-P_1/D)$, together with the one-dimensional
constituents already identified, give every GT path line over $F$.

The certificate chooses the first nonzero column of each rank-one
projector and divides it by its first nonzero entry.  The resulting
transition matrix is invertible over $F$.  The script checks that all
conjugated $P_i$ matrices agree exactly and that $P_1$ is diagonal with
the specified entries $D$ or $0$.  This fixes a reproducible rational
normalization and path order for this construction.  Under (N-25),
the associated matrix units are $F$-linear combinations of the
certified word basis, by solving the invertible 114-dimensional
linear system over $F$.  Thus they give a full GT matrix-unit basis
inside $\mathcal B_{4,F}$.

\paragraph{The remaining canonical-basis gate, precisely.}
GCT VIII, subsection 3.1, printed pp.~19--20, describes the GT path
construction under its multiplicity-free condition C1'.  Formula
(N-32) verifies that condition through rank four for the full algebra.
But VIII, subsection 3.3, printed pp.~20--22, additionally requires
the order induced from nonstandard duality, the scaling exponents
$a_T$, and unique global minimum-degree lifts.  These are its
Conjectures 3.1 and 3.2, not consequences of the split algebra.
Its subsection 5.2 computes a single distinguished rank-four element;
it supplies no complete rank-four table.

There is also a basic normalization issue that the degree comparison
alone cannot remove.  Suppose an initial GT vector $c'_T$ is replaced
by $u_Tc'_T$, with $u_T\in\mathbb Q^\times$.  Such changes arise, for
example, from a rational diagonal change of basis in an irreducible
module, with the dual factors transformed inversely.  Every lattice
generated in VIII's construction remains the same, because $u_T$ is
a unit in all the coefficient rings.  The residue label changes to
$u_Tb_T$, and its fibre is the scalar multiple by $u_T$ of the old
fibre.  For every candidate $y$, the ambient degree tuple satisfies
$d(u_Ty)=d(y)$.  Consequently, if the unique minimizing lift exists,
it changes to $u_Te_T$; the optimization in $a_T$ has the same
equivariance.  Thus a uniquely normalized table needs specified
initial constant normalizations or a further normalization rule.
The degree comparison by itself does not select them.

The present matrices establish neither that VIII's eventual $E(4)$
belongs to $\mathbb Q(D)$ nor that its signed factors lie in
$\mathbb N[q+q^{-1}]$.  They do show that an algebraic coefficient
extension cannot be justified merely by the need to split the
rank-four algebra or its GT tower.  Matrix units give structure
constants zero and one for their own $F$-basis, but this is no
dimension-ring assertion for VIII's designated basis and no integral
$A$-basis assertion, which N-009 excludes.

\paragraph{Literature status and verification.}
The nine generic simple models and generic dimension are known from
GCT IV.  GCT VII, subsection 7.2 (Example 2), explicitly states
splitting over its $\mathbb Q(q)$ in the paragraph following the
74-term relation figures, printed p.~40; its explicit simple-module
formulas follow in subsection 7.2.1.  Blasiak's \emph{Representation
theory of the nonstandard Hecke algebra}, Theorem 5.1 and section 6,
prove multiplicity-free branching and seminormal-basis results for
the nonstandard Temperley--Lieb quotient.  That quotient has generic
rank 89 at rank four, and it must not replace the full rank-114 algebra.
The contribution here is the explicit $\mathbb Q(D)$ descent, exact
full-algebra rank certificate, and complete rational GT path data;
the classification itself is not claimed as new.

The file \nolinkurl{rational_gt_rank4.py} contains the entire exact
construction and writes \nolinkurl{rational_gt_rank4_verified.json}.
It checks (N-28), symmetric/alternating restriction, quotient invariance,
the finite-field rank, all projectors, all path ranks, and all GT
intertwining identities.  The plethysm branch independently read the complete proof and script,
replayed the full construction, and checked the additional scalar and
rank-three parameter identities above.  Its whole-proof review is
\nolinkurl{nonstandard_plethysm_review_N010.md}, with the separate
additional-check receipt
\nolinkurl{rational_gt_rank4_plethysm_additional.json}.  The two
inherited inputs (generic dimension and $F$-form) retain that status.
